\chapter{Rời Ghế Nhà Trường, Đời Không Chấm Điểm Công Khai}
\markboth{Rời Ghế Nhà Trường, Đời Không Chấm Điểm Công Khai}{Leaving School, Entering an Ungraded Life}

\section{Ra Trường Xong, Không Ai Phát Đề Cương Cho Em Nữa}
\begin{truyen}{Ra Trường Xong, Không Ai Phát Đề Cương Cho Em Nữa}{After Graduation, No One Gives You an Outline}
\chuhoa{N}{gày còn đi học, em luôn biết kỳ thi ở đâu, môn nào quan trọng, và mốc nào phải chạy tới.}
\textit{(When you were in school, you always knew where the exam was, which subject mattered, and which deadline was coming.)}

Ra trường rồi, nhiều người hoảng vì không còn ai vẽ sẵn đường cho mình.
\textit{(After graduation, many people panic because no one draws the road for them.)}

Không có thời khóa biểu không có nghĩa là em tự do hơn ngay.
\textit{(No timetable does not mean instant freedom.)}

Nhiều khi nó chỉ có nghĩa là em phải tự cầm lái.
\textit{(It often only means that you must drive yourself.)}

Rất nhiều người trẻ không thua vì kém.
\textit{(Many young people do not fail because they are weak.)}

Họ thua vì quen sống trong đường ray quá lâu.
\textit{(They fail because they lived on rails for too long.)}

\end{truyen}

\begin{baihoc}
Bước vào đời là bước vào nơi không còn bài mẫu. Muốn đi được, em phải tập tự đặt mục tiêu và tự giữ nhịp cho mình.
\textit{(Entering adult life means entering a place with no model answer. To move well, you must set your own goals and keep your own rhythm.)}
\end{baihoc}

\begin{ghinhoanh}
\textbf{Short English to remember:}

Entering adult life means entering a place with no model answer.

To move well, you must set your own goals and keep your own rhythm.

\end{ghinhoanh}

\begin{tuhocnhanh}
\textbf{Cau hoi tu hoc / Self-study questions}

1. Van de chinh la gi? \textit{What is the main problem?}

2. Bai hoc thuc te la gi? \textit{What is the practical lesson?}

3. Mot cau tieng Anh em muon nho la cau nao? \textit{Which English sentence do you want to remember?}
\end{tuhocnhanh}
\ngancach

\section{Chiếc Đồng Phục Cởi Ra, Áp Lực Không Biến Mất}
\begin{truyen}{Chiếc Đồng Phục Cởi Ra, Áp Lực Không Biến Mất}{The Uniform Is Gone, But Pressure Stays}
\chuhoa{N}{hiều học sinh tưởng ra khỏi trường là thoát áp lực.}
\textit{(Many students think leaving school means escaping pressure.)}

Sự thật là áp lực chỉ đổi áo.
\textit{(The truth is that pressure only changes clothes.)}

Từ điểm số sang tiền bạc, từ kỳ thi sang công việc, từ kỳ vọng của thầy cô sang kỳ vọng của đời.
\textit{(It moves from scores to money, from exams to work, from teachers’ expectations to life’s expectations.)}

Nếu em không học cách chịu áp lực từ lúc còn trẻ, em sẽ bị đời dạy lại bằng học phí đắt hơn nhiều.
\textit{(If you do not learn how to carry pressure while young, life will teach it again at a much higher cost.)}

Trưởng thành không làm gánh nặng biến mất.
\textit{(Growing up does not make burdens disappear.)}

Nó chỉ yêu cầu em vững hơn khi mang nó.
\textit{(It only asks you to become steadier while carrying them.)}

\end{truyen}

\begin{baihoc}
Đừng ảo tưởng rằng ra trường là hết mệt. Điều quan trọng là học cách mệt mà vẫn không gãy.
\textit{(Do not imagine that life after school is free of strain. What matters is learning how to be tired without breaking.)}
\end{baihoc}

\begin{ghinhoanh}
\textbf{Short English to remember:}

Do not imagine that life after school is free of strain.

What matters is learning how to be tired without breaking.

\end{ghinhoanh}

\begin{tuhocnhanh}
\textbf{Cau hoi tu hoc / Self-study questions}

1. Van de chinh la gi? \textit{What is the main problem?}

2. Bai hoc thuc te la gi? \textit{What is the practical lesson?}

3. Mot cau tieng Anh em muon nho la cau nao? \textit{Which English sentence do you want to remember?}
\end{tuhocnhanh}
\ngancach

\section{Không Còn Chuông Vào Lớp, Cũng Không Còn Chuông Cứu Em}
\begin{truyen}{Không Còn Chuông Vào Lớp, Cũng Không Còn Chuông Cứu Em}{There Is No Bell to Save You Anymore}
\chuhoa{T}{rong trường học, mỗi tiết học đều có điểm bắt đầu và kết thúc rõ ràng.}
\textit{(In school, every class has a clear beginning and ending.)}

Ngoài đời, nhiều việc không tự dừng lại chỉ vì em đã mệt.
\textit{(Outside school, many things do not stop simply because you are tired.)}

Một công việc dở không có tiếng chuông cứu.
\textit{(A bad job has no bell to rescue you.)}

Một khoản nợ cũng không tự hết giờ.
\textit{(A debt does not expire because time is up.)}

Vì thế người trẻ phải học kỹ năng tự chặn mình, tự nghỉ đúng lúc, và tự quay lại đúng giờ.
\textit{(That is why young adults must learn to stop themselves, rest at the right time, and return at the right time.)}

\end{truyen}

\begin{baihoc}
Ra đời đòi hỏi tự quản nhiều hơn đi học. Ai không tự quản được thời gian và sức lực thì rất dễ bị cuốn trôi.
\textit{(Adult life demands more self-management than school. Those who cannot manage time and energy are easily swept away.)}
\end{baihoc}

\begin{ghinhoanh}
\textbf{Short English to remember:}

Adult life demands more self-management than school.

Those who cannot manage time and energy are easily swept away.

\end{ghinhoanh}

\begin{tuhocnhanh}
\textbf{Cau hoi tu hoc / Self-study questions}

1. Van de chinh la gi? \textit{What is the main problem?}

2. Bai hoc thuc te la gi? \textit{What is the practical lesson?}

3. Mot cau tieng Anh em muon nho la cau nao? \textit{Which English sentence do you want to remember?}
\end{tuhocnhanh}
\ngancach

\section{Bằng Cấp Là Vé Vào Cửa, Không Phải Ghế Ngồi Sẵn}
\begin{truyen}{Bằng Cấp Là Vé Vào Cửa, Không Phải Ghế Ngồi Sẵn}{A Degree Opens the Door, Not the Seat}
\chuhoa{C}{ó người cầm bằng tốt nghiệp và tưởng từ đây mình đương nhiên phải có chỗ đứng tốt.}
\textit{(Some people hold a degree and think that from now on a good place is automatically theirs.)}

Nhưng bằng cấp chỉ giúp em được nhìn thấy.
\textit{(But a degree only helps you be seen.)}

Nó không tự chứng minh em làm được việc.
\textit{(It does not prove that you can do the work.)}

Ngoài đời, chỗ đứng thường phải được dựng bằng việc làm thật, cách làm việc thật, và thái độ thật.
\textit{(In real life, your place is usually built by real work, real discipline, and real attitude.)}

Cầm bằng mà kiêu quá sớm là một kiểu ngây thơ rất phổ biến của tuổi trẻ.
\textit{(Holding a degree and becoming proud too early is a very common youth mistake.)}

\end{truyen}

\begin{baihoc}
Hãy tôn trọng bằng cấp, nhưng đừng thần thánh hóa nó. Đời còn hỏi thêm rất nhiều câu sau tấm bằng.
\textit{(Respect a degree, but do not worship it. Life asks many more questions after the certificate.)}
\end{baihoc}

\begin{ghinhoanh}
\textbf{Short English to remember:}

Respect a degree, but do not worship it.

Life asks many more questions after the certificate.

\end{ghinhoanh}

\begin{tuhocnhanh}
\textbf{Cau hoi tu hoc / Self-study questions}

1. Van de chinh la gi? \textit{What is the main problem?}

2. Bai hoc thuc te la gi? \textit{What is the practical lesson?}

3. Mot cau tieng Anh em muon nho la cau nao? \textit{Which English sentence do you want to remember?}
\end{tuhocnhanh}
\ngancach

\section{Ngày Đầu Không Còn Ai Gọi Em Là Học Sinh Giỏi}
\begin{truyen}{Ngày Đầu Không Còn Ai Gọi Em Là Học Sinh Giỏi}{On Day One, No One Calls You the Top Student}
\chuhoa{C}{ó những người từng rất sáng trong trường học nhưng bước vào môi trường mới lại thấy mình bình thường hẳn.}
\textit{(Some people were brilliant in school but enter a new environment and suddenly feel very ordinary.)}

Điều đó đau, nhưng cần.
\textit{(That hurts, but it is necessary.)}

Nó kéo cái tôi từ trên mây xuống mặt đất.
\textit{(It brings the ego down from the clouds to the ground.)}

Ở nơi mới, em không còn sống bằng danh hiệu cũ.
\textit{(In a new place, you cannot live on old titles.)}

Em phải bắt đầu lại bằng giá trị hiện tại.
\textit{(You must begin again with present value.)}

Bắt đầu lại từ số không không hạ nhục em.
\textit{(Starting again from zero does not humiliate you.)}

Nó chỉ nói rằng sân chơi đã đổi.
\textit{(It only means the arena has changed.)}

\end{truyen}

\begin{baihoc}
Người trẻ cần học cách rời bỏ vinh quang cũ. Ai ôm mãi quá khứ giỏi giang thường khó lớn tiếp.
\textit{(Young people must learn to leave old glory behind. Those who cling to yesterday’s success often struggle to grow.)}
\end{baihoc}

\begin{ghinhoanh}
\textbf{Short English to remember:}

Young people must learn to leave old glory behind.

Those who cling to yesterday’s success often struggle to grow.

\end{ghinhoanh}

\begin{tuhocnhanh}
\textbf{Cau hoi tu hoc / Self-study questions}

1. Van de chinh la gi? \textit{What is the main problem?}

2. Bai hoc thuc te la gi? \textit{What is the practical lesson?}

3. Mot cau tieng Anh em muon nho la cau nao? \textit{Which English sentence do you want to remember?}
\end{tuhocnhanh}
\ngancach

\section{Không Ai Nhắc Em Học Nữa}
\begin{truyen}{Không Ai Nhắc Em Học Nữa}{No One Reminds You to Learn Anymore}
\chuhoa{R}{a truong roi, khong con ai giao bai tap hay nhac em phai boi them ky nang.}
\textit{(After graduation, no one gives homework or reminds you to build new skills.)}

Nhiều người tưởng việc học sẽ tự dừng khi thời khóa biểu dừng lại.
\textit{(Many people think learning ends when the timetable ends.)}

Nhưng người lớn đi nhanh thường là người tự học khi không ai ép.
\textit{(But adults who move well are usually the ones who keep learning without pressure.)}

Ai chờ người khác nhắc mới học sẽ chậm rất lâu sau khi ra trường.
\textit{(Those who wait for reminders will stay slow for a long time.)}

\end{truyen}

\begin{baihoc}
Ra đời là lúc việc học chuyển từ nghĩa vụ sang trách nhiệm tự thân.
\textit{(Adult life turns learning from obligation into personal responsibility.)}
\end{baihoc}

\begin{ghinhoanh}
\textbf{Short English to remember:}

Adult life turns learning from obligation into personal responsibility.

\end{ghinhoanh}

\begin{tuhocnhanh}
\textbf{Cau hoi tu hoc / Self-study questions}

1. Van de chinh la gi? \textit{What is the main problem?}

2. Bai hoc thuc te la gi? \textit{What is the practical lesson?}

3. Mot cau tieng Anh em muon nho la cau nao? \textit{Which English sentence do you want to remember?}
\end{tuhocnhanh}
\ngancach

\section{Thứ Hai Đầu Tiên Không Có Lớp}
\begin{truyen}{Thứ Hai Đầu Tiên Không Có Lớp}{The First Monday with No Class}
\chuhoa{B}{uổi sáng đầu tuần đầu tiên sau khi rời trường làm nhiều người thấy hụt hẫng rất lạ.}
\textit{(The first Monday morning after leaving school feels strangely empty for many people.)}

Họ tưởng không có lịch học sẽ là nhẹ tênh và thoải mái.
\textit{(They think having no class schedule will feel light and free.)}

Nhưng khoảng trống không có cấu trúc rất dễ biến thành loay hoay và trì hoãn.
\textit{(But empty time without structure easily becomes confusion and delay.)}

Từ đó người trẻ mới hiểu tự lập bắt đầu từ việc tự chia ngày cho mình.
\textit{(Then young people understand that independence begins by shaping their own day.)}

\end{truyen}

\begin{baihoc}
Sự trống không không tự sinh ra tự do. Nó đòi em biết tổ chức mình.
\textit{(Emptiness does not automatically create freedom. It demands self-organization.)}
\end{baihoc}

\begin{ghinhoanh}
\textbf{Short English to remember:}

Emptiness does not automatically create freedom.

It demands self-organization.

\end{ghinhoanh}

\begin{tuhocnhanh}
\textbf{Cau hoi tu hoc / Self-study questions}

1. Van de chinh la gi? \textit{What is the main problem?}

2. Bai hoc thuc te la gi? \textit{What is the practical lesson?}

3. Mot cau tieng Anh em muon nho la cau nao? \textit{Which English sentence do you want to remember?}
\end{tuhocnhanh}
\ngancach

\section{Quá Nhiều Ngã Rẽ}
\begin{truyen}{Quá Nhiều Ngã Rẽ}{Too Many Crossroads}
\chuhoa{R}{ời trường xong, có người đứng trước quá nhiều lựa chọn nên không biết bước theo hướng nào.}
\textit{(After school, some people face too many options and do not know which way to take.)}

Họ tưởng nhiều lựa chọn đồng nghĩa nhiều thuận lợi.
\textit{(They think more choices always mean more advantage.)}

Nhưng nếu không hiểu mình, nhiều ngã rẽ chỉ làm đầu óc ồn hơn.
\textit{(But without self-understanding, many options only make the mind louder.)}

Người trưởng thành không chỉ chọn nhiều.
\textit{(Mature people do not only choose more.)}

Họ học cách bỏ bớt để đi chắc hơn.
\textit{(They learn to cut options and walk more steadily.)}

\end{truyen}

\begin{baihoc}
Không có đề cương sẵn nghĩa là em phải tự chọn câu hỏi quan trọng cho đời mình.
\textit{(No fixed outline means you must choose the important questions for your own life.)}
\end{baihoc}

\begin{ghinhoanh}
\textbf{Short English to remember:}

No fixed outline means you must choose the important questions for your own life.

\end{ghinhoanh}

\begin{tuhocnhanh}
\textbf{Cau hoi tu hoc / Self-study questions}

1. Van de chinh la gi? \textit{What is the main problem?}

2. Bai hoc thuc te la gi? \textit{What is the practical lesson?}

3. Mot cau tieng Anh em muon nho la cau nao? \textit{Which English sentence do you want to remember?}
\end{tuhocnhanh}
\ngancach

\section{Danh Hiệu Cũ Không Theo Em Mãi}
\begin{truyen}{Danh Hiệu Cũ Không Theo Em Mãi}{Old Titles Do Not Follow You Forever}
\chuhoa{C}{ó người từng sống rất lâu trong lời khen học sinh giỏi nên ra đời mới thấy hụt khi không ai còn gọi mình như vậy.}
\textit{(Some lived a long time inside praise as a top student, then felt empty when no one used that title anymore.)}

Họ tưởng tên gọi cũ sẽ tiếp tục bảo vệ giá trị hiện tại.
\textit{(They think old labels will keep protecting present value.)}

Ngoài đời, hôm nay em làm được gì quan trọng hơn em từng được gọi là gì.
\textit{(In real life, what you can do today matters more than what you were once called.)}

Rời danh hiệu cũ đau nhưng nó giúp em bước vào giá trị thật hơn.
\textit{(Leaving old titles hurts, but it leads into a more real value.)}

\end{truyen}

\begin{baihoc}
Muốn lớn tiếp, người trẻ phải đủ can đảm sống mà không dựa mãi vào hào quang cũ.
\textit{(To keep growing, young people need the courage to stop leaning on old glory.)}
\end{baihoc}

\begin{ghinhoanh}
\textbf{Short English to remember:}

To keep growing, young people need the courage to stop leaning on old glory.

\end{ghinhoanh}

\begin{tuhocnhanh}
\textbf{Cau hoi tu hoc / Self-study questions}

1. Van de chinh la gi? \textit{What is the main problem?}

2. Bai hoc thuc te la gi? \textit{What is the practical lesson?}

3. Mot cau tieng Anh em muon nho la cau nao? \textit{Which English sentence do you want to remember?}
\end{tuhocnhanh}
\ngancach

\section{Không Ai Chấm Nỗ Lực Công Khai}
\begin{truyen}{Không Ai Chấm Nỗ Lực Công Khai}{No One Publicly Grades Your Effort}
\chuhoa{Đ}{i học, cố gắng thường được điểm số phản ánh phần nào.}
\textit{(In school, effort is often reflected in grades at least partly.)}

Nhiều người tưởng ra đời rồi cố là sẽ có ai đó công nhận rõ ràng như thế.
\textit{(Many people think adult life will recognize effort just as clearly.)}

Nhưng đời ít chấm công khai.
\textit{(But life rarely grades openly.)}

Có khi em cố rất lâu mà không nghe một lời khen nào.
\textit{(Sometimes you work hard for a long time without hearing praise.)}

Lúc đó điều giữ em đi tiếp không còn là điểm, mà là nhận thức mình đang xây gì.
\textit{(At that point, what keeps you moving is no longer a grade, but knowing what you are building.)}

\end{truyen}

\begin{baihoc}
Trưởng thành là làm nhiều việc đúng mà không cần bảng điểm xác nhận mỗi lần.
\textit{(Maturity means doing many right things without needing a score each time.)}
\end{baihoc}

\begin{ghinhoanh}
\textbf{Short English to remember:}

Maturity means doing many right things without needing a score each time.

\end{ghinhoanh}

\begin{tuhocnhanh}
\textbf{Cau hoi tu hoc / Self-study questions}

1. Van de chinh la gi? \textit{What is the main problem?}

2. Bai hoc thuc te la gi? \textit{What is the practical lesson?}

3. Mot cau tieng Anh em muon nho la cau nao? \textit{Which English sentence do you want to remember?}
\end{tuhocnhanh}
